%%%==============EX_1============%%%
\begin{ex}%[2H5V2-6]
Trong không gian $Ox y z$, cho hai điểm $A(2; 1;-3)$, $B(-3; 2; 1)$. Gọi $(d)$ là đường thẳng đi qua $M(1; 2; 3)$ sao cho tổng khoảng cách từ $A$ đến $(d)$ và từ $B$ đến $(d)$ là lớn nhất. Khi đó phương trình đường thẳng $(d)$ là
	\choice
	{$\dfrac{x-1}{-5}=y-2=\dfrac{z}{4}$}
	{$\dfrac{x-1}{-3}=\dfrac{y-2}{2}=\dfrac{z-3}{1}$}
	{\True $\dfrac{x-1}{1}=\dfrac{y-2}{13}=\dfrac{z-3}{-2}$}
	{$\dfrac{x-1}{-3}=\dfrac{y-2}{2}=\dfrac{z-3}{2}$}
	\loigiai{
		Ta có $\mathrm{d}(A,(d)) \leq AM$, $\mathrm{d}(B,(d)) \leq BM$.\\
		$\mathrm{d}(A,(d))+\mathrm{d}(B,(d))$ đạt giá trị lớn nhất khi và chỉ khi $\heva{&AM\perp(d) \\
			&BM\perp(d).}$\\		
$\overrightarrow{MA}=(1;-1;-6)$ và $\overrightarrow{MB}=(-4; 0;-2)$.\\	
Khi đó $(d)$ có véc-tơ  chỉ phương $\vec{u}=\left[ \overrightarrow{MA}, \overrightarrow{MB}\right] =(2; 26;-4) \Rightarrow \overrightarrow{u}_1=(1; 13;-2)$.\\
Phương trình $(d)\colon \dfrac{x-1}{1}=\dfrac{y-2}{13}=\dfrac{z-3}{-2}$.
	}
\end{ex}
%%%==============EX_2============%%%
\begin{ex}%[2H5V2-6]
Trong không gian $Ox y z$, cho bốn điểm $A(3; 0; 0)$, $B(0; 2; 0)$, $C(0; 0; 6)$ và $D(1; 1; 1)$. Gọi $\Delta$ là đường thẳng đi qua $D$ và thỏa mãn tổng khoảng cách từ các điểm $A$, $B$, $C$ đến $\Delta$ là lớn nhất. Hỏi $\Delta$ đi qua điểm nào trong các điểm dưới đây?
	\choice
	{\True  $M(5; 7; 3)$ }
	{$M(-1;-2; 1)$}
	{$M(3; 4; 3)$}
	{$M(7; 13; 5)$} 
	\loigiai{
		Gọi $(P)$ là mặt phẳng đi qua $3$ điểm $A$, $B$, $C$.
		\begin{center}
\begin{tikzpicture}[>=stealth,line join=round,line cap=round,font=\footnotesize,scale=.7]
	\def\r{5}% Bán kính cầu
	\def\h{1.75}%Chiều cao IH
	\def\g{10}%Góc tiếp xúc cầu và Ellipse
	\pgfmathsetmacro{\am}{sqrt((\r)^2-(\h)^2)}
	\pgfmathsetmacro{\a}{\am *sec(\g)}
	\pgfmathsetmacro{\b}{\a/4}
	\pgfmathsetmacro{\bm}{\b *sin(\g)}
	\pgfmathsetmacro{\gm}{asin(\h/\r)}
	\path
	(0:0) coordinate (I)
	(270:\h-\bm) coordinate (O) node[below]{$D$}
	($(O)+(-165:2*\r)$) coordinate (P1)
	($(O)+(160:1.5*\r)$) coordinate (P2)
	($(P1)+(0:3.5*\r)$) coordinate (P4)
	($(P2)+(0:3.5*\r)$) coordinate (P3)
	($(O)+(-30:0.7*\r)$) coordinate (C)
($(O)+(-150:0.5*\r)$) coordinate (A)	
($(O)+(50:0.5*\r)$) coordinate (B)
($(O)+(-60:0.7*\r)$) coordinate (K)
	($(O)+(-60:0.3*\r)$) coordinate (J)
	($(O)+(120:0.5*\r)$) coordinate (H)
	($(O)+(120:1.5*\r)$) coordinate (d) node[below]{$\Delta$}
	($(O)+(-60:\r)$) coordinate (d');
		\draw (P1)--(P2)--(P3)--(P4)--cycle 
		(O)--(d) (H)--(A)--(B)--(C) (K)--(d');
		\draw [dashed] (d')--(H) (B)--(J) (C)--(K);
\foreach \x/\g in {A/-170,B/20,C/0,J/-110,K/-130,H/-170}
\draw[fill=black] 	(\x) circle (.5pt)
($(\g:.3)+(\x)$) node {$\x$};			
\end{tikzpicture}
		\end{center}
Vì ba điểm $A(3; 0; 0)$, $B(0; 2; 0)$, $C(0; 0; 6)$ nằm trên ba trục tọa độ nên ta có phương trình của $(P)\colon \dfrac{x}{3}+\dfrac{y}{2}+\dfrac{z}{6}=1\Leftrightarrow 2x+3y+z-6=0$.\\		
		Ta thấy $D(1; 1; 1) \in(P)$.\\		
		Gọi $H$, $J$, $K$ lần lượt là hình chiếu của $A$, $B$, $C$ lên $\Delta$.\\		
Khi đó $AH\leq AD$, $BJ\leq BD$, $CK\leq CD$.\\
 Suy ra tổng khoảng cách từ các điểm $A$, $B$, $C$ đến $\Delta$ là lớn nhất bằng $AD+BD+CD$ hay $\Delta$ là đường thẳng đi qua $D$ và vuông góc với $(P)$.\\
	Suy ra phương trình $\Delta\colon\heva{&x=1+2t \\
		&y=1+3t \\
		&z=1+t.}$\\
		Khi đó $\Delta$ đi qua $M(5; 7; 3)$.
	}
\end{ex}
%%%==============EX_3============%%%
\begin{ex}%[2H5V2-6]
	Trong không gian $Ox y z$, cho hai đường thẳng $\Delta\colon \dfrac{x}{1}=\dfrac{y}{1}=\dfrac{z-1}{1}$ và $\break \Delta'\colon\dfrac{x-1}{1}=\dfrac{y}{2}=\dfrac{z}{1}$. Xét điểm $M$ thay đổi. Gọi $a$, $b$ lần lượt là khoảng cách từ $M$ đến $\Delta$ và $\Delta'$. Biểu thức $a^2+2b^2$ đạt giá trị nhỏ nhất khi và chỉ khi $M\equiv M_0\left(x_0, y_0, z_0\right)$. Khi đó giá trị $x_0+y_0$ bằng
	\choice
	{$\dfrac{4}{3}$}
	{$0$}
	{\True $\dfrac{2}{3}$}
	{$\sqrt{2}$}
	\loigiai{
		Giả sử $PQ$ là đoạn vuông góc chung của hai đường thẳng $\Delta$ và $\Delta'$.\\	
		Với $P\in \Delta$ và $Q\in \Delta'$ có $P(0,0,1)$ và $Q(1,0,0)$. Gọi $E$, $F$ lần lượt là hình chiếu của $M$ trên đường thẳng $\Delta$ và $\Delta'$.\\		
		Khi đó ta có $a+b=M E+M F \geq E F \geq PQ$ hay $a+b \geq \sqrt{2}$.\\	
		Nhận thấy $\dfrac{1}{3}(a-2 b)^2=\left(a^2+2b^2\right)-\left(\dfrac{2}{3} a^2+\dfrac{4}{3} a b+\dfrac{2}{3} b^2\right)=\left(a^2+2b^2\right)-\dfrac{2}{3}(a+b)^2$.\\
		Do đó ta có $\dfrac{1}{3}(a-2 b)^2 \geq 0 \Leftrightarrow a^2+2b^2 \geq \dfrac{2}{3}(a+b)^2$ hay $a^2+2b^2 \geq \dfrac{4}{3}$.\\
		Vậy giá trị nhỏ nhất của $a^2+2b^2$ bằng $\dfrac{4}{3}$ đạt được khi dấu bằng xảy ra khi $\heva{&M \in E F \\ &E F \equiv PQ \\ &a=2 b}$ hay $\overrightarrow{MP}=-2\overrightarrow{MQ}$ hay $M\left(\dfrac{2}{3}, 0, \dfrac{1}{3}\right)$.
	}
\end{ex}
%%%==============EX_4============%%%
\begin{ex}%[2H5V2-6]
	Trong không gian $Ox y z$, cho điểm $A(1; 1; 2)$ và mặt phẳng $\break (P)\colon (m-1) x+y+m z-1=0$, với $m$ là tham số. Biết khoảng cách từ điểm $A$ đến mặt phẳng $(P)$ lớn nhất. Khẳng định đúng trong bốn khẳng định dưới đây là
	\choice
	{\True $2< m < 6$}
	{$m > 6$}
	{$-2< m < 2$}
	{$-6< m < 2$}
	\loigiai{
		\textbf{Cách 1:}\\
		Ta có $\mathrm{d}(A;(P))=\dfrac{|m-1+1+2m-1|}{\sqrt{(m-1)^2+1+m^2}}=\sqrt{\dfrac{(3m-1)^2}{2\left(m^2-m+1\right)}}$.\\		
		Xét $f(m)=\dfrac{(3m-1)^2}{2\left(m^2-m+1\right)} \Rightarrow f'(m)=\dfrac{(5-m)(3m-1)}{2\left(m^2-m+1\right)^2}=0\Leftrightarrow\hoac{&m=\dfrac{1}{3} \\
			&m=5.}$
		\begin{center}
			\begin{tikzpicture}[scale=1, font=\footnotesize, line join=round, line cap=round, >=stealth]
				\tkzTabInit[lgt=1,espcl=3,deltacl=0.5]
				{$t$/1.2,$f'(t)$/0.8,$f(t)$/2.5}
				{$-\infty$,$\dfrac{1}{3}$,$5$,$+\infty$}
				\tkzTabLine{,-,$0$,+,$0$,-,} %d là không xác định
				\tkzTabVar{+/$\dfrac{9}{2}$,-/$0$,+/$\dfrac{14}{3}$,-/$\dfrac{9}{2}$}
			\end{tikzpicture}		
		\end{center}
		Vậy $\max \mathrm{d}(A;(P))=\sqrt{\dfrac{14}{3}}$ khi $m=5\in(2; 6)$.\\		
		\textbf{Cách 2:}\\		
		Ta đi tìm đối tượng cố định của mặt phẳng $(P)$.\\		
		$(P)\colon (m-1) x+y+m z-1=0 \Leftrightarrow(x+z) m-x+y-1=0$.\\		
		Với mọi $m \in \mathbb{R}$ mặt phẳng $(P)$ luôn đi qua giao tuyến của hai mặt phẳng $\heva{&x+z=0\\
			&-x+y-1=0}$ tức luôn đi qua đường thẳng $(d)\colon\heva{&x=t \\
			&y=1+t \\
			&z=-t.} $\\		
		Gọi $H(t; 1+t;-t) \in(d) \Rightarrow \overrightarrow{AH}=(t-1; t;-t-2)$. Để khoảng cách từ $A$ đến $(P)$ lớn nhất thì $AH\perp(P) \Rightarrow \overrightarrow{AH}$ cùng phương với VTPT của $(P)$ là $\vec{n}_P=(m-1; 1; m)$, suy ra		
		\[\dfrac{t-1}{m-1}=\dfrac{t}{1}=\dfrac{-t-2}{m} \Leftrightarrow\heva{&m t-t=t-1\\
			&m t=-t-2} \Leftrightarrow\heva{&t=-\dfrac{1}{3} \\
			&m=5\in(2; 6).}\]
	}
\end{ex}
%%%==============EX_5============%%%

\begin{ex}%[2H5V2-6]
	Trong không gian $Ox y z$, cho điểm $A(0; 3;-2)$. Xét đường thẳng $d$ thay đổi, song song với trục $Oz$ và cách trục $Oz$ một khoảng bằng $2$. Khi khoảng cách từ $A$ đến $d$ lớn nhất, $d$ đi qua điểm nào dưới đây?
	\choice
	{$M(0; 8;-5)$}
	{$N(0; 2;-5)$}
	{\True $P(0;-2;-5)$}
	{$Q(-2; 0;-3)$}
	\loigiai{
		\begin{center}
			\begin{tikzpicture}[scale=1, font=\footnotesize, line join=round, line cap=round, >=stealth]
				\def\x{3}
				\pgfmathsetmacro{\y}{\x/4}
				\def\h{3}
				\coordinate[label=below left:$H$] (O) at (0,0);
				\coordinate[] (A') at (\x,0);
				\coordinate[label=left:$B$] (B) at (-\x,0);
				\coordinate[] (S) at (0,\h);
				\path ($(B)+(0,4)$) coordinate (d)node[left]{$d$}
				($(B)+(0,-2)$) coordinate (d')
				($(O)+(0,4)$) coordinate (z)  node[left]{$Oz$}
				($(O)+(0,-2)$) coordinate (z')
				($(A')+(0,-2)$) coordinate (e')
				($(A')+(0,4)$) coordinate (e)
				($(A')+(1.5,0)$) coordinate (A) node[below]{$A(0;3;-2)$};
				\draw (A') arc (0:-180:{\x} and {\y});
				\draw[dashed] (A') arc (0:180:{\x} and {\y}) (O)--(B);
				\coordinate[] (M) at ($(O)+(-120:3cm and .75cm)$);
				\draw[dashed](A')--(B) (z)--(z') ;
				\draw (d)--(d') (e)--(e') (A')--(A);
				\foreach \x/\g in {B/180}
				\draw[fill=black] 	(\x) circle (.5pt)
				($(\g:.3)+(\x)$) node {$\x$};	
			\end{tikzpicture}
		\end{center}
		Do đường thẳng $d\parallel Oz$ nên $d$ nằm trên mặt trụ có trục là $Oz$ và bán kính trụ là $R=2$.\\
		Gọi $H$ là hình chiếu của $A$ trên trục $Oz$, suy ra tọa độ $H(0; 0;-2)$.\\	
		Do đó $\mathrm{d}(A, Oz)=AH=3$.\\
		Gọi $B$ là điểm thuộc đường thẳng $AH$ sao cho $\overrightarrow{AH}=\dfrac{3}{5} \overrightarrow{AB}\Rightarrow B(0;-2;-2)$.\\		
		Vậy $\mathrm{d}(A, d)_{\max}=5\Leftrightarrow d$ là đường thẳng đi qua $B$ và song song với $Oz$.\\
		Phương trình tham số của $d\colon\heva{&x=0\\&y=-2\\&z=-2+t.}$\\		
		Kết luận $d$ đi qua điểm $P(0;-2;-5)$.
	}
\end{ex}
%%%==============EX_6============%%%
\begin{ex}%[2H5C2-7]
	Trong không gian $Ox y z$, đường thẳng $\Delta$ đi qua điểm $M(3; 1; 1)$, nằm trong mặt phẳng $(\alpha)\colon x+y-z-3=0$ và tạo với đường thẳng $d\colon\heva{&x=1\\
		&y=4+3t \\
		&z=-3-2t
	}$ một góc nhỏ nhất thì phương trình của $\Delta$ là
	\choice
	{$\heva{&x=1\\
			&y=-t' \\
			&z=2t'}$}
	{\True $\heva{&x=8+5t' \\
			&y=-3-4t' \\
			&z=2+t'}$}
	{$\heva{&x=1+2t' \\
			&y=1-t' \\
			&z=3-2t'}$}
	{$\heva{&x=1+5t' \\
			&y=1-4t' \\
			&z=3+2t'}$}
	\loigiai{
		\begin{center}
			\begin{tikzpicture}[>=stealth,line join=round,line cap=round,font=\footnotesize,scale=.5]
				\def\r{7}% Bán kính cầu
				\def\h{1.75}%Chiều cao IH
				\def\g{10}%Góc tiếp xúc cầu và Ellipse
				\pgfmathsetmacro{\am}{sqrt((\r)^2-(\h)^2)}
				\pgfmathsetmacro{\a}{\am *sec(\g)}
				\pgfmathsetmacro{\b}{\a/4}
				\pgfmathsetmacro{\bm}{\b *sin(\g)}
				\pgfmathsetmacro{\gm}{asin(\h/\r)}
				\path
				(0:0) coordinate (I)
				(270:\h-\bm) coordinate (O)
				($(O)+(-165:2*\r)$) coordinate (P1)
				($(O)+(160:1.5*\r)$) coordinate (P2)
				($(P1)+(0:3.5*\r)$) coordinate (P4)
				($(P2)+(0:3.5*\r)$) coordinate (P3)
				($(O)+(-130:0.3*\r)$) coordinate (M)
				($(M)+(50:\r)$) coordinate (N)	
					($(M)+(50:1.3*\r)$) coordinate (d_1)	node[below]{$d_1$}
					($(M)+(-130:0.7*\r)$) coordinate (m)
			($(M)+(-130:0.4*\r)$) coordinate (m1)
				($(M)+(10:\r)$) coordinate (H) 		
			($(M)+(10:1.3*\r)$) coordinate (de) 			node[below]{$\Delta$}
			($(M)+(-170:0.6\r)$) coordinate (de1)
				($(M)+(25:0.7*\r)$) coordinate (K)
		($(O)+(160:0.7*\r)$) coordinate (d')
		($(d')+(50:1*\r)$) coordinate (d) node[below]{$d$};		
				\draw (P1)--(P2)--(P3)--(P4)--cycle 
				(M)--(d_1) (de)--(de1) (m1)--(m) (N)--(H) (d)--(d');
				\draw [dashed] (M)--(K)--(H) (N)--(K) (M)--(m1);
				\foreach \x/\g in {M/-40,H/-60,K/30,N/130}
				\draw[fill=black] 	(\x) circle (.5pt)
				($(\g:.3)+(\x)$) node {$\x$};			
			\end{tikzpicture}
		\end{center}
		Đường thẳng $d$ có véc-tơ  chỉ phương là $\vec{u}=(0; 3;-2)$.\\		
		Mặt phẳng $(\alpha)$ có véc-tơ  pháp tuyến là $\vec{n}=(1; 1;-1)$.\\
		Vì $\vec{u} \cdot \vec{n}=0\cdot 1+3\cdot 1+(-2) \cdot(-1)=5\neq 0$ nên $d$ cắt $(\alpha)$.\\
		Gọi $d_1$ là đường thẳng đi qua $M$ và $d_1 \parallel d$, suy ra $d_1$ có phương trình $\heva{&x=3\\
			&y=1+3t \\
			&z=1-2t.}$\\
		Lấy $N(3; 4;-1) \in d_1$. Gọi $K$, $H$ lần lượt là hình chiếu vuông góc của $N$ trên mặt phẳng $(\alpha)$ và đường thẳng $\Delta$.\\		
		Ta có $\widehat{(d, \Delta)}=\widehat{NMH}$ và $\sin \widehat{NMH}=\dfrac{NH}{MN} \geq \dfrac{NK}{MN}$.\\		
		Do vậy $\widehat{(d, \Delta)}$ nhỏ nhất khi $K\equiv H$ hay $\Delta$ là đường thẳng $MK$.\\
		Đường thẳng $NK$ có phương trình $\heva{&x=3+t \\
			&y=4+t \\
			&z=-1-t.}$\\
		Tọa độ điểm $K$ ứng với $t$ là nghiệm của phương trình
		$$(3+t)+(4+t)-(-1-t)-3=0 \Leftrightarrow t=-\dfrac{5}{3} \text {. Suy ra}\, K\left(\dfrac{4}{3}; \dfrac{7}{3}; \dfrac{2}{3}\right).	$$
	}
\end{ex}
%%%==============EX_7============%%%
\begin{ex}%[2H5C3-3]
	Trong không gian $Ox y z$, cho điểm $A(0; 1;-2)$, mặt phẳng $(P)\colon x+y+z+1=0$ và mặt cầu $(S)\colon x^2+y^2+z^2-2x-4y-7=0$. Gọi $\Delta$ là đường thẳng đi qua $A$ và $\Delta$ nằm trong mặt phẳng $(P)$ và cắt mặt cầu $(S)$ tại hai điểm $B$, $C$ sao cho tam giác $IBC$ có diện tích lớn nhất, với $I$ là tâm của mặt cầu $(S)$. Phương trình của đường thẳng $\Delta$ là
	\choice
	{$\heva{&x=t \\
			&y=1\\
			&z=-2-t}$}
	{$\heva{&x=t \\
			&y=1+t \\
			&z=-2+t}$}
	{\True $\heva{&x=t \\
			&y=1-t \\
			&z=-2}$}
	{$\heva{&x=t \\
			&y=1+t \\
			&z=-2}$}
	\loigiai{
		\begin{center}
			\begin{tikzpicture}[>=stealth,line join=round,line cap=round,font=\footnotesize,scale=.5]
				\def\b{4}
				\def\a{5}
				\coordinate (I) at (0,0);
				\coordinate (A) at (0,-3);
				\draw[dashed] (4,-3) arc (0:180:4 and 0.8);
				\draw (4,-3) arc (0:-180:4 and 0.8);
				\draw (5,0) arc (0:360:5);
				\coordinate (B) at ($(A)+({4*cos(-60)},{0.8*sin(-60)})$);
				\path 
	 ($(B)!2!(A)$) coordinate (C);
				\coordinate (A') at (-8,-4);
				\coordinate (B') at (-4,-4);
				\coordinate (C') at (-5,-1);
				\draw (C')--(A')--(B');
				\draw[dashed] (I)--(A) (I)--(C)--(B)--(I);
				\draw (-8,-3.8)node[right]{$P$};
				\coordinate (r') at ($(M)!.5!(A)$);
				\coordinate (r) at ($(M)!.5!(I)$);
				\coordinate (d) at ($(A)!.5!(I)$);
				\foreach \i/\g in {I/90, A/-90, B/-45,C/150}
				\fill (\i) circle (1.5pt)+(\g:3mm) node {$\i$};
	\draw pic[draw,angle radius=2mm]{right angle=B--A--I};
			\end{tikzpicture}
		\end{center}
		$(S)$ có tâm $I(1; 2; 0)$ và bán kính $R=\sqrt{1^2+2^2+7}=2\sqrt{3}$.\\
	Ta có	$\overrightarrow{AI}=(1; 1; 2) \Rightarrow AI=\sqrt{6} < R.$ Suy ra $A$ nằm trong mặt cầu $(S)$ và $A$ nằm trên dây cung $BC$ \qquad{(1)}.\\
		$S_{\triangle IBC}=\dfrac{1}{2} I B \cdot I C \cdot \sin \widehat{B I C}=\dfrac{R^2}{2} \sin \widehat{B I C} \leq \dfrac{R^2}{2}$.\\
		Khi đó diện tích $\triangle IBC$ đạt giá trị lớn nhất là
		$\dfrac{R^2}{2} \Leftrightarrow \sin \widehat{B I C}=1 \Rightarrow \widehat{B I C}=90^{\circ} \Rightarrow \triangle IBC$ vuông cân tại $I \Rightarrow BC=I C \sqrt{2}=R \sqrt{2}=2 \sqrt{6}$.\\
		Gọi $J$ là trung điểm của $BC$. Ta có $I J \perp BC$ và $I J=\dfrac{BC}{2}=\sqrt{6}$ \qquad{(2)}.\\		
		$\triangle AIJ$ vuông tại $J\Rightarrow AI\geq IJ$, kết hợp thêm với (1) và (2) ta có $IJ=AI\Rightarrow A\equiv J\Rightarrow A$ là trung điểm của $BC$ và $I A \perp BC$.\\
		$(P)$ có véc-tơ  pháp tuyến $\vec{n}_{(P)}=(1; 1; 1)$ có giá vuông góc với $\Delta$.\\
		Vậy $\Delta$ nhận $\vec{u}=\left[\vec{n}_{(P)}, \overrightarrow{AI}\right]=(1;-1; 0)$ làm véc-tơ  chỉ phương và đi qua $A(0; 1;-2)$\\
		$\Rightarrow \Delta\colon \heva{&x=t \\
			&y=1-t\\
			&z=-2.}$
	}
\end{ex}

%%%==============EX_8============%%%
\begin{ex}%[2H5C3-3]
	Trong không gian $Ox y z$, cho điểm $E(2; 1; 3)$, mặt phẳng $(P)\colon 2x+2y-z-3=0$ và mặt cầu $(S)\colon(x-3)^2+(y-2)^2+(z-5)^2=36$. Gọi $\Delta$ là đường thẳng đi qua $E$, nằm trong $(P)$ và cắt $(S)$ tại hai điểm có khoảng cách nhỏ nhất. Phương trình của $\Delta$ là
	\choice
	{$\heva{&x=2+9t \\
			&y=1+9t \\
			&z=3+8t}$}
	{$\heva{&x=2-5t \\
			&y=1+3t \\
			&z=3}$}
	{\True $\heva{&x=2+t \\
			&y=1-t \\
			&z=3}$}
	{$\heva{&x=2+4t \\
			&y=1+3t \\
			&z=3-3t}$}	
	\loigiai{
		\begin{center}			
			\begin{tikzpicture}[line width=0.6pt,line cap=round,line join=round,>=stealth,scale=0.6]
				\def\r{4}
	\path
	(0:0) coordinate (I)
	($(I)+(-150:\r)$) coordinate (A)
	($(I)+(-30:\r)$) coordinate (B)
	($(I)+(-90:0.5*\r)$) coordinate (E)	
	($(I)+(-60:\r)$) coordinate (D)
	 ($(D)!3!(E)$) coordinate (C');	
%	\path  	[name path=c1] (I) circle (4);
%	\path  	[name path=de] (D)--(E);
%	\path  	[name intersections={of=c1 and de,by={C}}]	
	%\path [name path=c1] (O_1)  let\p1=($(O_1)-(A)$) in circle ({veclen(\x1,\y1});
	%\path [name path=c2] (O_2) circle (\r);
	
%	\path [name intersections={of=c1 and de,by={D,C}}];
		\path[name path=c2] let\p2=($(I)-(B)$) in (I) circle ({veclen(\x2,\y2)});
	\path[name path=d2] (C')--(D);
	\path[name intersections={of=c2 and d2,by={C,D}}];
\path	 ($(C)!0.5!(D)$) coordinate (H);	
%	($(I)+(-90:0.5*\r)$) coordinate (C)
				\draw (I) circle(4);
				\draw (H)--(I)--(E) (C)--(D) (A)--(B);
				%	\tkzMarkRightAngle[size=0.13](I,H,A)
\foreach \x/\g in {A/-170,B/-20,C/170,D/-70,E/-130,H/-120,I/0}
\draw[fill=black] 	(\x) circle (1pt)
($(\g:.3)+(\x)$) node {$\x$};				
			\end{tikzpicture}
		\end{center}
		Mặt cầu $(S)$ có tâm $I(3; 2; 5)$ và bán kính $R=6$.\\
		$IE=\sqrt{1^2+1^2+2^2}=\sqrt{6} < R\Rightarrow$ điểm $E$ nằm trong mặt cầu $(S)$.\\
		Gọi $H$ là hình chiếu của $I$ trên mặt phẳng $(P)$, $A$ và $B$ là hai giao điểm của $\Delta$ với $(S)$.\\
		Khi đó, $AB$ nhỏ nhất $\Leftrightarrow AB \perp O E$, mà $AB \perp IH$ nên $AB \perp(H I E) \Rightarrow AB \perp I E$.\\
		Suy ra $\overrightarrow{u}_{\Delta}=\left[\overrightarrow{n}_P; \overrightarrow{EI}\right]=(5;-5; 0)=5(1;-1; 0)$.\\
		Vậy phương trình của $\Delta$ là $\heva{&x=2+t \\
			&y=1-t \\
			&z=3.}$
	}
\end{ex}
%%%==============EX_9============%%%

\begin{ex}%[2H5C3-3]
	Trong không gian $Ox y z$, cho điểm $M(-3; 3;-3)$ thuộc mặt phẳng $\break (\alpha)\colon 2x-2y+z+15=0$ và mặt cầu $(S)\colon(x-2)^2+(y-3)^2+(z-5)^2=100$. Đường thẳng $\Delta$ qua $M$, nằm trên mặt phẳng $(\alpha)$ cắt $(S)$ tại $A$, $B$ sao cho độ dài $AB$ lớn nhất. Viết phương trình đường thẳng $\Delta$.
	\choice
	{$\dfrac{x+3}{1}=\dfrac{y-3}{1}=\dfrac{z+3}{3}$}
	{\True $\dfrac{x+3}{1}=\dfrac{y-3}{4}=\dfrac{z+3}{6}$}
	{$\dfrac{x+3}{16}=\dfrac{y-3}{11}=\dfrac{z+3}{-10}$}
	{$\dfrac{x+3}{5}=\dfrac{y-3}{1}=\dfrac{z+3}{8}$}
		\loigiai{
Ta có mặt cầu $(S)$ có tâm $I(2; 3; 5)$, bán kính $R=10$.\\
$\mathrm{d}(I,(\alpha))=\dfrac{|2\cdot2-2\cdot3+5+15|}{\sqrt{2^2+(-2)^2+1^2}}=6< R\Rightarrow(\alpha) \cap(S)=C(H; r)$, với $H$ là hình chiếu của $I$ lên $(\alpha)$.\\
Gọi $\Delta_1$ là đường thẳng qua $I$ và vuông góc với $(\alpha) \Rightarrow \Delta_1$ có véc-tơ  chỉ phương là $\break \overrightarrow{u}_{\Delta_1}=(2;-2; 1)$.\\
			Phương trình tham số của $ \Delta_1\colon\heva{&x=2+2t \\
				&y=3-2t \\
				&z=5+t.}$\\	
			Tọa độ $H$ là nghiệm của hệ $\heva{&x=2+2t \\
				&y=3-2t \\
				&z=5+t \\
				&2x-2y+z+15=0} \quad \Rightarrow\heva{&x=-2\\
				&y=7\\
				&z=3} \Rightarrow H(-2; 7; 3)$.\\		
			Ta có $AB$ có độ dài lớn nhất khi và chỉ khi $AB$ là đường kính của $(C)\Leftrightarrow \Delta \equiv MH$.\\		
			Đường thẳng $MH$ đi qua $M(-3; 3;-3)$ và có véc-tơ  chỉ phương $\overrightarrow{MH}=(1; 4; 6)$.\\
			Suy ra phương trình $\Delta\colon \dfrac{x+3}{1}=\dfrac{y-3}{4}=\dfrac{z+3}{6}$.
		}
	\end{ex}


%%%==============EX_10============%%%
\begin{ex}%[2H5C3-3]
	Trong không gian $Ox y z$, cho điểm $E(2; 1; 3)$, mặt phẳng $(P)\colon 2x+2y-z-3=0$ và mặt cầu $(S)\colon (x-3)^2+(y-2)^2+(z-5)^2=36$. Gọi $\Delta$ là đường thẳng đi qua $E$, nằm trong mặt phẳng $(P)$ và cắt $(S)$ tại hai điểm có khoảng cách nhỏ nhất. Phương trình của $\Delta$ là
	\choice
	{$\heva{&x=2+9t \\
			&y=1+9t \\
			&z=3+8t}$}
	{$\heva{&x=2-5t \\
			&y=1+3t \\
			&z=3}$}
	{\True $\heva{&x=2+t \\
			&y=1-t \\
			&z=3}$}
	{$\heva{&x=2+4t \\
			&y=1+3t \\
			&z=3-3t}$}
	\loigiai{
		Ta có tâm và bán kính mặt cầu $(S)$ là $I(3; 2; 5)$; $R=6$.\\
		$IE=\sqrt{1+1+4}=\sqrt{6} < R$.\\
		Gọi $\Delta$ là đường thẳng đi qua $E$.\\
		Gọi $H$ là hình chiếu vuông góc của $I$ lên $\Delta$.\\
		Dây cung càng nhỏ khi khoảng cách từ tâm tới đường thẳng $\Delta$ càng lớn.\\
		Ta có $\mathrm{d}(I, \Delta)=IH\leq IE$.\\
		Vậy dây cung nhỏ nhất khi đường thẳng $\Delta$ vuông góc với $\overrightarrow{IE}=(-1;-1;-2)$.\\
		Dựa vào các đáp án ta thấy trong các véc-tơ  chỉ phương $\overrightarrow{u}_1=(9; 9; 8)$,  $\overrightarrow{u}_2=(-5; 3; 0)$, $\overrightarrow{u}_3=(1;-1; 0)$,		
		$\overrightarrow{u}_4=(4; 3;-3)$
		thì chỉ có $\overrightarrow{u}_3 \cdot \overrightarrow{IE}=0$.\\		
		\textbf{Nhận xét}: ta hoàn toàn có thể viết được phương trình đường thẳng $\Delta$ bằng cách viết phương trình mặt phẳng $(Q)$ đi qua $E$ nhận $\overrightarrow{IE}=(-1;-1;-2)$ làm một véc-tơ  pháp tuyến, khi đó $\Delta=(P) \cap(Q)$.
	}
\end{ex}
%%%==============EX_11============%%%
\begin{ex}%[2H5C3-4]
	Trong không gian $Ox y z$, cho mặt cầu $(S)\colon x^2+y^2+z^2-2x+4y+2z-3=0$ và mặt phẳng $(P)\colon 2x-y+2z-14=0$. Điểm $M$ thay đổi trên $(S)$, điểm $N$ thay đổi trên $(P)$. Độ dài nhỏ nhất của $MN$ bằng
	\choice
	{\True $1$}
	{$2$}
	{$\dfrac{1}{2}$}
	{$\dfrac{3}{2}$}
	\loigiai{
		\begin{center}		
				\begin{tikzpicture}[scale=.7, font=\footnotesize, line join=round, line cap=round, >=stealth]
					\def\R{3}
					\coordinate[label=left:$ $] (I) at (0,0);
					\coordinate[label=right:$ $] (B) at (-\R,0);
					\coordinate[label=right:$ $] (A) at (\R,0);
					\draw (I) circle (\R) (A) arc (0:-180: {\R} and {\R/3});
					\draw[dashed] (A) arc (0:180: {\R} and {\R/3});
					\path ($(I)+(-90:\R)$) coordinate (M)
					($(I)+(-90:1.3*\R)$) coordinate (H);
				%	\coordinate[] (M) at ($(I)!1!-45:(A)$);
				%	\coordinate (N) at ($(I)!1!-135:(A)$);				
	\def\r{3}% Bán kính cầu
	\path (2,-3.9) coordinate (N)
	%(0:0) coordinate (I)
	($(H)+(-165:2*\r)$) coordinate (P)
	($(H)+(160:1.5*\r)$) coordinate (P2)
	($(P)+(0:3.5*\r)$) coordinate (P4)
	($(P2)+(0:3.5*\r)$) coordinate (P3);
\draw (P)--(P4) (P)--(P2)  (M)--(N)--(H)--(M);
\draw [dashed] (I)--(M);		
\foreach \x/\g in {M/-160,H/180,N/0,I/0,P/30}
\draw[fill=black] 	(\x) circle (1pt)
($(\g:.3)+(\x)$) node {$\x$};	

	\draw pic[draw,,angle radius=6mm]{angle=P4--P--P2};%Theo chiều dương, tùy chọn double,fill=yellow!50,-stealth,dashed
	%	\coordinate[label=right:$H$] (H) at ($ (M)!0.5!(N) $);
	%	\draw (M) arc (0:-180:1.75cm and 0.3cm);
	%	\draw[dashed] (M) arc (0:180:1.75cm and 0.3cm);
	%	\draw[dashed] (I);
	%		\foreach \diem in {A,B,I,M,N}	\fill (\diem)circle(1.5pt);	
				\end{tikzpicture}			
		\end{center}
		Mặt cầu $(S)$ có tâm $I(1;-2;-1)$, bán kính $R=3$; $\mathrm{d}(I;(P))=4> R\Rightarrow$ mặt cầu $(S)$ và mặt phẳng $(P)$ không có điểm chung.\\		
		Dựng $IH\perp(P)$, $(H\in(P))$.\\
		Ta có $MN$ nhỏ nhất khi $M$ là giao điểm của đoạn $IH$ với $(S)$ và $N\equiv H$.\\		
		Phương trình đường thẳng $IH\colon\heva{&x=1+2t \\
			&y=-2-t \\
			&z=-1+2t}$ ; $(t \in \mathbb{R})$.\\
		Điểm $M(1+2t;-2-t;-1+2t) \in(S)$ nên $(x-1)^2+(y+2)^2+(z+1)^2=9$\\
		$$\Leftrightarrow(2t)^2+(-t)^2+(2t)^2=9\Leftrightarrow t=\pm 1.$$
		Khi đó $M_1(3;-3; 1)$, $M_2(-1;-1;-3)$.\\
		Thử lại $\mathrm{d}\left(M_1;(P)\right)=1$; $\mathrm{d}\left(M_2;(P)\right)=7 > IH=4$ (loại).\\
		Vậy $MN_{\min}=MH=1$ khi $M(3;-3; 1)$; $N\left(\dfrac{11}{3};-\dfrac{10}{3}; \dfrac{5}{3}\right)$.
	}
\end{ex}
%%%==============EX_12============%%%
\begin{ex}%[2H5C3-4]
	Trong không gian $Ox y z$, cho các điểm $A(1; 0; 0)$, $B(3; 2; 0)$, $C(-1; 2; 4)$. Gọi $M$ là điểm thay đổi sao cho đường thẳng $MA$, $MB$, $MC$ hợp với mặt phẳng $(ABC)$ các góc bằng nhau; $N$ là điểm thay đổi nằm trên mặt cầu $(S)\colon(x-3)^2+(y-2)^2+(z-3)^2=\dfrac{1}{2}$. Tính giá trị nhỏ nhất của độ dài đoạn $MN$.
	\choice
	{$\dfrac{3\sqrt{2}}{2}$}
	{$\sqrt{2}$}
	{\True $\dfrac{\sqrt{2}}{2}$}
	{$\sqrt{5}$}
	\loigiai{
\begin{center}		
	\begin{tikzpicture}[scale=.7, font=\footnotesize, line join=round, line cap=round, >=stealth]
	\def\a{4}
	\def\h{4}
	\path 	(0:0) coordinate (A)
			++(0:\a) coordinate (B)
			++(-150:4*\a/5) coordinate (C)
			($(C)!1/2!(B)$) coordinate (H)
			($(H)+(90:\h)$) coordinate (S);
	\draw[thick] (C)--(A) (C)--(B)
			(A)--(S)	(B)--(S)--(H)	(C)--(S);
	\draw[dashed,thick] (A)--(B);
	\foreach \x / \goc in 		{A/180,B/0,C/-135,H/-35,S/90}
			\fill (\x) circle (1.5pt)
			($(\x)+(\goc:3mm)$) node {$\x$};
	
	\draw pic[draw,angle radius=2mm]{right angle=B--H--S};%Theo chiều dương
\end{tikzpicture}			
\end{center}	
Ta có $\overrightarrow{A B}=(2 ; 2 ; 0)$, $\overrightarrow{A C}=(-2 ; 2 ; 4) \Rightarrow \overrightarrow{A B} \cdot \overrightarrow{A C}=0 \Rightarrow \triangle ABC$ vuông tại $A$.\\
Gọi $H$ là hình chiếu vuông góc của điểm $M$ trên mặt phẳng $(A B C)$. Ta có
$$
\begin{aligned}
	& (MA, (ABC))=(MA, HA)=\widehat{MAH} \\
	& (MB, (ABC))=(MB, HB)=\widehat{MBH} \\
	& (MC, (ABC))=(MC, HC)=\widehat{MCH}.
\end{aligned}
$$
 Theo giả thiết $\widehat{MAH}=\widehat{MBH}=\widehat{MCH} \Rightarrow \triangle MAH=\triangle MBH=\triangle MCH$ (g.c.g).\\
Do đó $HA=HB=HC$ nên $H$ là tâm đường tròn ngoại tiếp $\triangle ABC$.\\
Suy ra $H$ là trung điểm của $BC\Rightarrow H\ (1; 2; 2)$.\\
Ta có $[\overrightarrow{AB}, \overrightarrow{AC}]=(8;-8; 8)$, chọn véc-tơ  chỉ phương của đường thẳng $MH$ là $\overrightarrow{u}_{MH}=(1;-1; 1)$. \\
Phương trình đường thẳng $MH$ có dạng $\heva{&x=1+t\\&y=2-t\\&z=2+t}$, $t\in \mathbb{R}$.\\
	Mặt cầu $(S)$ có tâm $I(3 ; 2 ; 3)$ và bán kính $R=\dfrac{\sqrt{2}}{2}$.\\
 Gọi $K(1+t; 2-t; 2+t)$ là hình chiếu vuông góc của điểm $I$ trên đường thẳng $MH$.\\
Ta có $\overrightarrow{IK}=(t-2;-t; t-1), \overrightarrow{u_{MH}}=(1;-1; 1)$.\\
Do $IK\perp MH$ nên $\overrightarrow{IK} \cdot \overrightarrow{u_{MH}}=0$, ta được $t=1$.\\
Khi đó $K\ (2; 1; 3)$ và $IK=\sqrt{2}$.\\
Do $IK > R$ nên đường thẳng $MH$ không cắt mặt cầu.\\
Ta có $MN\geq \mathrm{d}(I, MH)-IN=IK-IN=\dfrac{\sqrt{2}}{2}$.\\
Vậy giá trị nhỏ nhất của độ dài đoạn $M N$ bằng $\dfrac{\sqrt{2}}{2}$.
	}
\end{ex}